\documentclass[10pt,a4paper]{article}
\usepackage[margin=20mm,headheight=14pt]{geometry}
\usepackage[T1]{fontenc}
\usepackage{lmodern}
\usepackage{microtype}
\usepackage{xcolor}
\usepackage{booktabs}
\usepackage{tabularx}
\usepackage{array}
\usepackage{longtable}
\usepackage{enumitem}
\usepackage{tikz}
\usetikzlibrary{arrows.meta,positioning,shapes.geometric}
\usepackage{fancyhdr}
\usepackage[hidelinks]{hyperref}
\usepackage{listings}
\usepackage{titlesec}

\definecolor{ehgreen}{HTML}{245B43}
\definecolor{ehmint}{HTML}{EAF2ED}
\definecolor{ehgray}{HTML}{56625B}
\definecolor{ehline}{HTML}{D9E1DB}
\hypersetup{colorlinks=true,linkcolor=ehgreen,urlcolor=ehgreen}
\pagestyle{fancy}
\fancyhf{}
\lhead{\small EasyHousing | Low-Level Design}
\rhead{\small Version 1.0}
\cfoot{\thepage}
\setlength{\parskip}{4pt}
\setlength{\parindent}{0pt}
\setlist[itemize]{leftmargin=17pt,itemsep=1pt,topsep=2pt}
\setlist[enumerate]{leftmargin=19pt,itemsep=1pt,topsep=2pt}
\titleformat{\section}{\Large\bfseries\color{ehgreen}}{\thesection}{0.7em}{}
\titleformat{\subsection}{\large\bfseries\color{ehgray}}{\thesubsection}{0.7em}{}
\newcommand{\code}[1]{\texttt{\detokenize{#1}}}
\newcommand{\endpoint}[1]{\path{#1}}
\newcommand{\field}[1]{\texttt{#1}}
\newcommand{\pagebreakdoc}{\clearpage}

\begin{document}
\begin{titlepage}
\centering
\vspace*{19mm}
{\color{ehgreen}\rule{\textwidth}{2pt}}\par
\vspace{12mm}
{\Huge\bfseries EasyHousing\par}
\vspace{4mm}
{\LARGE Low-Level Design\par}
\vspace{5mm}
{\large\color{ehgray}Service contracts, workflows, data model, and implementation details}\par
\vspace{13mm}
{\color{ehgreen}\rule{0.62\textwidth}{1pt}}\par
\vspace{16mm}
\begin{tabular}{rl}
\textbf{Document type:} & Low-Level Design (LLD)\\
\textbf{System:} & EasyHousing property discovery and listing platform\\
\textbf{Backend:} & ASP.NET Core / .NET 8, Entity Framework Core\\
\textbf{Frontend:} & Angular 22, TypeScript, RxJS, server-rendering support\\
\textbf{Data:} & SQL Server (LocalDB in current development configuration)\\
\textbf{Revision:} & 1.0\\
\textbf{Prepared:} & November 2026\\
\end{tabular}
\vfill
\begin{minipage}{0.88\textwidth}
\small\color{ehgray}
This document is based on the source and configuration present in the repository.
All example ports and development behavior refer to the current local configuration.
No production deployment, external payment flow, automated API test suite, or
unimplemented infrastructure is implied.
\end{minipage}
\vspace{8mm}
\end{titlepage}

\tableofcontents
\pagebreakdoc

\section{Purpose, Scope, and Implementation Baseline}
This Low-Level Design describes how EasyHousing's implemented frontend and backend
parts cooperate. It is intended for developers who need to trace a UI action to an
HTTP route, service method, data record, and response. The companion HLD describes
system boundaries and quality attributes.

\subsection{Repository layout}
\begin{verbatim}
EasyHousingSolution/
|-- EasyHousingSolution.sln
|-- backend/
|   |-- backendd.sln
|   |-- docs/                 (this document and HLD source)
|   `-- services/
|       |-- ApiGateway/       (src/Program.cs)
|       |-- IdentityService/  (src/Controllers, Data, Domain, Services)
|       |-- UserService/      (src/Controllers, Data, Domain)
|       |-- PropertyService/  (src/Controllers, Data, Domain, Services)
|       `-- CartService/      (src/Controllers, Data, Domain, Services)
`-- frontend/
    `-- src/app/
        |-- core/             (configuration, models, guards, interceptors, clients)
        `-- features/          (auth, buyer, seller, admin, shared)
\end{verbatim}

Each backend project keeps its \code{.csproj}, settings, and launch profile at the
service root; application code is under \code{src/}. Each project is an independent
ASP.NET Core web application targeting .NET 8. The frontend is a standalone Angular
application with route-level feature components.

\subsection{Runtime and configuration}
\begin{tabularx}{\textwidth}{>{\bfseries}p{27mm} p{20mm} X}
\toprule
Process & Port & Important configuration\\
\midrule
Angular dev server & 4200 & \code{API_BASE_URL=http://localhost:5200/api}; local session in browser storage.\\
API Gateway & 5200 & Downstream URLs in \code{Services:*}; CORS origin list.\\
Identity Service & 5201 & Identity database, JWT issuer/audience/signing key, User Service URL, internal key.\\
User Service & 5202 & User database, same JWT validation settings, same internal key.\\
Property Service & 5203 & Property database, JWT settings, User Service URL and internal key.\\
Cart Service & 5204 & Cart database, JWT settings, Property Service URL.\\
\bottomrule
\end{tabularx}
Development launch profiles set \code{ASPNETCORE_ENVIRONMENT=Development}, which
enables Swagger. The service defaults use LocalDB and development keys; production
must replace these values with secure, environment-specific configuration.

\subsection{Cross-cutting request processing}
The gateway maps the first segment following \code{/api/} to a downstream base URL
and streams request bodies and content headers. It copies status and response headers
back to the caller. It does not own application authorization and currently contains
no configured retry, circuit-breaker, or health-check behavior. Each service configures
controllers, SQL Server EF Core, CORS, Swagger, and JWT validation as appropriate.
API exception middleware maps selected service failures to JSON error responses.

\pagebreakdoc
\section{API Gateway and HTTP Contract Conventions}
\subsection{Gateway routing table}
\begin{tabularx}{\textwidth}{>{\ttfamily}p{42mm} >{\ttfamily}p{40mm} X}
\toprule
Public prefix & Downstream key & Service path behavior\\
\midrule
/api/identity/* & IdentityService & Path and query are forwarded unchanged.\\
/api/locations/* & UserService & Public state/city endpoints.\\
/api/properties/* & PropertyService & Search, seller, admin, image and detail endpoints.\\
/api/cart/* & CartService & Buyer saved-home endpoints.\\
\bottomrule
\end{tabularx}
An unknown prefix returns 404. Gateway configuration keys are loaded from
\code{Services:IdentityService}, \code{Services:UserService},
\code{Services:PropertyService}, and \code{Services:CartService}. The API gateway
also serves a development Swagger UI and proxies each service's
\code{/swagger/v1/swagger.json}.

\subsection{Conventions}
\begin{itemize}
  \item JSON DTO properties are serialized using ASP.NET Core web JSON defaults
  (camel case); TypeScript models use camel-case fields.
  \item Authenticated browser calls receive \code{Authorization: Bearer <JWT>} from
  the Angular HTTP interceptor when the URL is under the configured API origin.
  \item Role-protected endpoints return the platform's standard ASP.NET Core
  authorization outcomes; controllers additionally return explicit 400, 404, 409,
  201, 200, or 204 outcomes where implemented.
  \item Internal User Service calls pass \code{X-Internal-Api-Key}. Do not expose this
  header or value to the browser.
  \item Request cancellation tokens are passed through to EF Core and downstream HTTP
  calls for the service operations inspected here.
\end{itemize}

\subsection{High-level route inventory}
\begin{longtable}{p{24mm} p{69mm} p{44mm}}
\toprule
Service & Route(s) & Access\\
\midrule
\endhead
Identity & \endpoint{POST /api/identity/register}; \endpoint{POST /api/identity/login} & Public\\
User & \endpoint{GET /api/locations/states}; \endpoint{GET /api/locations/states/\{stateId\}/cities} & Public\\
User & \endpoint{GET /api/internal/locations/resolve}; \endpoint{/api/internal/users*} & Internal API key\\
Property & \endpoint{GET /api/properties/search}; \endpoint{GET /api/properties/\{id\}} & Public active/verified data\\
Property & \endpoint{GET /api/properties/compare}; \endpoint{GET /api/properties/\{id\}/contact} & Buyer\\
Property & \endpoint{POST /api/properties}; \endpoint{PUT /api/properties/\{id\}}; image upload; \endpoint{GET /api/properties/mine} & Seller\\
Property & Pending, owner, region, verify, deactivate admin routes & Admin\\
Cart & \endpoint{GET /api/cart}; \endpoint{POST /api/cart}; \endpoint{DELETE /api/cart/\{cartId\}} & Buyer\\
\bottomrule
\end{longtable}

\pagebreakdoc
\section{Identity Service: Accounts, Registration, and Login}
\subsection{Responsibilities and persistence}
Identity Service owns \code{IdentityAccount}: \field{IdentityAccountId} (primary key),
\field{UserName}, \field{PasswordHash}, \field{UserType}, and \field{UserId}. EF Core
enforces a unique index on username. Username and role lengths are capped at 25 and
15 characters; the hash column is capped at 500 characters. \field{UserId} is the
profile identifier in User Service, not a cross-database foreign key.

\subsection{Registration contract}
\begin{tabularx}{\textwidth}{>{\bfseries}p{36mm} X}
\toprule
Field & Validation / meaning\\
\midrule
\field{userName} & Required, 1--25 characters; trimmed before account creation.\\
\field{password} & Required, 8--25 characters; hashed by ASP.NET Core PasswordHasher.\\
\field{userType} & Exactly Buyer or Seller; Admin registration is not accepted.\\
\field{firstName}, \field{lastName} & First name required, each at most 25 characters.\\
\field{emailId} & Required valid email, at most 50 characters.\\
\field{phoneNo} & Required exactly ten digits.\\
\field{dateOfBirth} & Date in the past; checked by controller.\\
\field{address}, \field{stateId}, \field{cityId} & Seller profile requires address and valid state/city selection.\\
\bottomrule
\end{tabularx}
The route first checks username existence and account type/date constraints. Identity
then sends a profile creation request to User Service. User Service additionally
checks duplicate username/email and, for sellers, verifies that the city belongs to
the selected state. On profile success, Identity hashes the password and saves the
account. A database update conflict triggers a compensating delete of the profile
and returns a username conflict response. This is a small synchronous compensation,
not a distributed transaction; email-uniqueness conflict behavior should be verified
against the exact exception path if the registration contract is changed.

\subsection{Login and token contract}
Login selects the account by trimmed username and verifies the submitted password
with the stored hash. Unknown account or incorrect password returns 401 with a
generic message. Successful login returns:
\begin{verbatim}
{
  "userName": "seller.name",
  "userType": "Seller",
  "profileId": 123,
  "token": "<signed JWT>"
}
\end{verbatim}
The JWT uses HMAC SHA-256 with configured signing key; claims include subject, name,
role, and custom \field{profileId}. Expiration is eight hours. Services validate
issuer, audience, signing key, lifetime, and a 30-second clock skew. The local admin
account is seeded at startup only if absent and uses a development password; replace
or disable it for any shared deployment.

\subsection{Identity error mapping}
Argument validation exceptions are returned as 400. A downstream/operation
\code{InvalidOperationException} is logged as a warning and returned as 502 with a
registration-availability message. Unexpected exceptions are logged and returned as
500. Explicit controller responses include duplicate username 409 and invalid login
401. Error JSON uses a \field{Message} property; frontend code may also accept
lower-case \field{message} from proxied components.

\pagebreakdoc
\section{User Service: Profiles and Locations}
\subsection{Owned tables}
\begin{longtable}{>{\bfseries}p{28mm} p{41mm} p{79mm}}
\toprule
Entity & Key fields & Constraints and relationships\\
\midrule
\endhead
\field{UserProfile} & \field{UserId} PK; username, role, names, date of birth, phone, email, optional address/state/city & Unique username and email; optional State and City references; seller location is mandatory at create time.\\
\field{State} & \field{StateId} PK; \field{StateName} & State name max 30; one-to-many City relationship.\\
\field{City} & \field{CityId} PK; \field{CityName}; \field{StateId} FK & City name max 50; referenced State delete is restricted.\\
\bottomrule
\end{longtable}
User profile first/last names are capped at 25 characters; email max 50, username
max 25. User Service database startup seeds Gujarat, Maharashtra, Karnataka and
Ahmedabad, Surat, Mumbai, Pune, Bengaluru. Seeding checks for existing rows before
adding; the initial catalog is deliberately small and is not an exhaustive
geographic reference dataset.

\subsection{Location endpoints}
\begin{longtable}{p{38mm} p{22mm} p{74mm}}
\toprule
Endpoint & Status & Behavior\\
\midrule
\endhead
\endpoint{GET /api/locations/states} & 200 & State IDs and names ordered by name.\\
\endpoint{GET /api/locations/states/\{stateId\}/cities} & 200 & City IDs, names, and state IDs for the selected state, ordered by city name.\\
\endpoint{GET /api/internal/locations/resolve?stateId=\&cityId=} & 200/401/404 & Internal-key check; confirms selected city is in selected state; returns normalized region.\\
\bottomrule
\end{longtable}
The property submission stores both state and city IDs and a resolved region string.
This allows filtering by region while preserving normalized identifiers in property
data. If a location no longer resolves, Property Service maps the not-found result
to a client validation error.

\subsection{Internal profile/contact endpoints}
\begin{longtable}{p{38mm} p{22mm} p{74mm}}
\toprule
Endpoint & Status & Behavior\\
\midrule
\endhead
\endpoint{POST /api/internal/users} & 200/400/401/409 & Internal key; creates buyer/seller profile and returns \field{userId}.\\
\endpoint{DELETE /api/internal/users/\{userId\}} & 204/401/404 & Internal key; used for registration compensation.\\
\endpoint{GET /api/internal/users/\{userId\}/contact} & 200/401/404 & Internal key; returns owner name, email, and phone.\\
\bottomrule
\end{longtable}
The internal-key comparison is ordinal and expects the configured
\code{InternalApiKey}. Unauthorized internal calls return 401. API exception
middleware logs unhandled errors and emits a generic JSON 500. Avoid logging
credentials, access tokens, or the internal key when diagnosing failures.

\subsection{Account/profile relationship}
The Identity account stores the profile ID returned by User Service. Registration
therefore creates the User Service row before the Identity row. User Service is the
source of truth for contact fields; Identity is the source of truth for credentials
and role. Property and cart services use profile IDs from the JWT to associate
records without direct database joins.

\pagebreakdoc
\section{Property Service: Listing Lifecycle and Query Rules}
\subsection{Property aggregate}
\field{PropertyListing} owns its \field{PropertyImage} collection.
\begin{longtable}{>{\bfseries}p{32mm} p{39mm} p{75mm}}
\toprule
Group & Fields & Rules / meaning\\
\midrule
\endhead
Identity & PropertyId, SellerId & Property ID is local PK; SellerId is a User Service profile ID.\\
Listing & PropertyName, PropertyType, PropertyOption, Description & Name max 50; type max 15; option Sell/Rent; description max 250.\\
Location & Address, Region, StateId, CityId, Landmark & Address max 250; landmark max 25; region is resolved from User Service.\\
Pricing & PriceRange, InitialDeposit & Decimal precision 18,2; price must be positive; deposit non-negative.\\
Lifecycle & Status, IsActive, DeactivationReason, CreatedAtUtc & New listing pending/inactive; timestamps initialized in UTC.\\
Images & ImageId, PropertyId, Image, ContentType & Child rows; cascade delete with property.\\
\bottomrule
\end{longtable}

\subsection{State transitions}
\begin{center}
\begin{tikzpicture}[
  state/.style={draw=ehgreen,fill=ehmint,rounded corners,minimum width=27mm,
    minimum height=10mm,align=center,font=\small},
  edge/.style={-{Stealth[length=2mm]},thick,draw=ehgreen}
]
\node[state] (pending) {Pending\\inactive};
\node[state,right=35mm of pending] (verified) {Verified\\active};
\node[state,below=18mm of verified] (deactivated) {Deactivated\\inactive};
\draw[edge] (pending) -- node[above,align=center,font=\scriptsize]{Admin verify} (verified);
\draw[edge] (verified) -- node[right,align=left,font=\scriptsize]{Admin deactivate\\with reason} (deactivated);
\draw[edge] (verified) to[bend left=28] node[above,align=center,font=\scriptsize]{Seller edit: reset to Pending/inactive} (pending);
\draw[edge] (deactivated) to[bend right=25] node[below,align=center,font=\scriptsize]{Seller edit: reset to Pending/inactive} (pending);
\end{tikzpicture}
\end{center}
Public search and detail endpoints require both \field{Status = Verified} and
\field{IsActive = true}. Seller \code{mine} has a verified/active filter when
\code{status=Verified}; otherwise it returns non-verified or inactive records.
Admin pending returns records whose status is Pending. Admin region/owner lookup
returns matching records and includes their image collections.

\subsection{Property validation and service integration}
Create and update accept a JSON \field{PropertySubmission}. Data annotations enforce
many lengths/ranges; controller validation also rejects whitespace values, unsupported
Sell/Rent values, invalid state/city IDs, non-positive price, negative deposit, and
descriptions over 250 characters. The User Directory Client adds
\code{X-Internal-Api-Key} and resolves the location before persistence.

Seller identity comes from the JWT \field{profileId} claim; it is not accepted from
the body. Update selects by both property ID and authenticated seller ID. Admin
verification changes lifecycle flags; deactivation requires a non-empty reason of
at most 250 characters. The current admin region/owner endpoints use their Admin
role attribute.

\subsection{Property endpoint detail}
\begin{longtable}{p{45mm} p{21mm} p{68mm}}
\toprule
Endpoint & Access & Contract summary\\
\midrule
\endhead
\endpoint{GET /api/properties/search} & Public & Optional region/type/option/sortBy; active verified listings; includes images.\\
\endpoint{GET /api/properties/\{id\}} & Public & Single active verified property with images; otherwise 404.\\
\endpoint{GET /api/properties/compare?id1=\&id2=} & Buyer & Two distinct positive IDs; returns two active verified records or 400/404.\\
\endpoint{GET /api/properties/\{id\}/contact} & Buyer & Resolves seller profile and requests owner contact from User Service.\\
\endpoint{POST /api/properties} & Seller & Creates Pending listing; 201 with ID and message.\\
\endpoint{PUT /api/properties/\{id\}} & Seller & Updates own listing, resets it to pending/inactive; 204.\\
\endpoint{GET /api/properties/mine} & Seller & Lists seller-owned verified or non-verified/inactive properties.\\
\endpoint{GET /api/properties/admin/pending} & Admin & Pending moderation queue ordered by creation time.\\
\endpoint{PUT /api/properties/admin/\{id\}/verify} & Admin & Activates and verifies; 204.\\
\endpoint{PUT /api/properties/admin/\{id\}/deactivate} & Admin & Deactivates with required reason; 204.\\
\bottomrule
\end{longtable}

\pagebreakdoc
\section{Property Service: Image Upload and Delivery}
\subsection{Upload contract}
The seller UI creates a \code{FormData} and appends every selected file under the
field name \code{files}. It posts to
\endpoint{POST /api/properties/\{propertyId\}/images} with the seller JWT. The
controller uses multipart model binding to receive \code{List<IFormFile>}; no
explicit multipart Content-Type should be set by the frontend because the browser
must supply the boundary.

\begin{tabularx}{\textwidth}{>{\bfseries}p{42mm} X}
\toprule
Rule & Current implementation\\
\midrule
Ownership & Property must exist and its \field{SellerId} must equal JWT \field{profileId}; otherwise 404.\\
Cardinality & At least one selected file; total image count after upload must not exceed six.\\
Per-file size & Greater than zero and at most 5,000,000 bytes.\\
MIME allowlist & \code{image/jpeg}, \code{image/png}, \code{image/webp}; this checks declared Content-Type.\\
Request size & Action allows up to 32,000,000 bytes.\\
Persistence & File bytes and Content-Type are written to Property Service's SQL database as child image rows.\\
Success & 200 with a message; seller UI reloads listing data after upload.\\
\bottomrule
\end{tabularx}

\subsection{Read path and browser rendering}
Search/details/seller/admin queries eagerly load the \field{Images} navigation.
Property Service serializes each image into \field{imageId}, Base64 \field{image},
and \field{contentType}. Cart's downstream Property DTO mirrors this shape, allowing
saved homes to render the same image. The frontend shared
\code{propertyImageUrl()} helper:
\begin{enumerate}
  \item Returns the placeholder when no image data is present.
  \item Preserves already-formed data, HTTP(S), protocol-relative, and non-Base64
  site-relative paths.
  \item Wraps Base64 content in a \code{data:<mime>;base64,<payload>} URL, preserving
  JPEG/PNG/WebP content types.
\end{enumerate}
Image-bearing views bind the result to an HTML \code{img src}. This matters because
Angular's CSS style sanitization and slash-prefixed Base64 payloads can otherwise
turn a valid image into a blocked style or bogus site path. Actual image elements
also preserve intrinsic image decoding behavior while CSS \code{object-fit: cover}
handles cropping.

\subsection{Operational and security considerations}
At maximum limits, a property may store 30 MB of raw image payload before database
and Base64/JSON overhead. Search can therefore return large responses, and six files
may exceed the action request limit depending on their combined sizes. The allowlist
validates the client-declared MIME type only; production hardening should inspect
file signatures, apply upload rate and request limits, consider image re-encoding,
and serve thumbnails from object storage or a dedicated media endpoint. Current
support stores bytes in SQL and does not implement image deletion, reordering,
thumbnails, virus scanning, or signed download URLs.

\subsection{Image failure triage}
\begin{enumerate}
  \item Check upload response status and exact multipart field name \code{files}.
  \item Confirm listing ownership, total image count, file length, and allowed MIME type.
  \item Read a seller/admin listing response and confirm \field{images} is non-empty
  and every item has \field{image} and \field{contentType}.
  \item Confirm the browser source starts with \code{data:image/<type>;base64,} for
  Base64 data, not \code{/9j/...} as a relative URL.
  \item Check the actual image's browser \field{naturalWidth} and \field{naturalHeight};
  a non-empty string alone does not prove successful image decoding.
\end{enumerate}

\pagebreakdoc
\section{Cart Service: Saved Homes}
\subsection{Persistence model}
\field{SavedProperty} has \field{CartId} primary key, \field{BuyerId},
\field{PropertyId}, and UTC \field{AddedAtUtc}. EF Core enforces a unique composite
index on buyer/property. IDs are local references; the Cart database does not
foreign-key into the Property or User database.

\subsection{Cart operations}
\begin{longtable}{p{42mm} p{20mm} p{68mm}}
\toprule
Endpoint & Access & Behavior\\
\midrule
\endhead
\endpoint{GET /api/cart} & Buyer & Selects current buyer's rows newest first; resolves each property over HTTP; returns nullable property DTO.\\
\endpoint{POST /api/cart} & Buyer & Validates positive property ID; resolves public property; rejects missing or duplicate; persists the reference; returns 201.\\
\endpoint{DELETE /api/cart/\{cartId\}} & Buyer & Deletes only a row owned by current buyer; 404 if absent; otherwise 204.\\
\bottomrule
\end{longtable}
Buyer identity is taken from JWT \field{profileId}. A property not found/active at
add-time is rejected. During list retrieval, a missing property is represented as a
null property, preserving the saved row so the buyer can remove it. Cart exception
middleware maps missing downstream connectivity to 502 and logs it; unexpected
failures become generic 500 responses.

\subsection{Dependency and consistency characteristics}
Cart reads are N+1 synchronous downstream requests: one SQL query for saved items,
followed by one property HTTP request per row. This is simple and returns current
property state/images but makes latency and availability depend on Property Service.
There is no bulk property endpoint or cache in the current implementation. If cart
size grows, prefer a bounded batch endpoint and explicit behavior for unavailable
properties; do not silently return stale copied property data without a consistency
policy.

\section{Angular Frontend: Routes, Services, and State}
\subsection{Route and role model}
\begin{tabularx}{\textwidth}{>{\ttfamily}p{29mm} p{30mm} X}
\toprule
Route & Access & Main component\\
\midrule
/search & Public & Search and filter verified listings.\\
/property/:id & Public & Listing details; buyer-only contact action.\\
/login, /register & Public & Login and buyer/seller registration.\\
/compare, /cart & Buyer & Compare two listings; manage saved homes.\\
/seller & Seller & Seller listing form, listing status and photo upload.\\
/admin & Admin & Pending verification and property moderation.\\
/error & Public & Shared error view; unknown paths redirect here.\\
\bottomrule
\end{tabularx}
Angular route guards improve navigation and UX, but service-side
\code{[Authorize(Roles=...)]} is the actual authorization boundary.

\subsection{Core client responsibilities}
\begin{itemize}
  \item \service{AuthService} calls login/register, persists the returned session and
  token in local storage, exposes the current session as an Angular signal, and clears
  it on logout.
  \item \service{PropertyService} centralizes property, compare, contact, cart,
  seller, and admin HTTP operations. Image upload uses \code{FormData}.
  \item \service{LocationService} requests state and state-specific city lists.
  \item \service{AuthInterceptor} adds Bearer credentials to requests matching the
  configured API base URL in a browser context; it does not add local storage tokens
  during server-side rendering.
  \item \service{ErrorInterceptor} routes network, server, and selected not-found
  errors to a shared page, with property search excluded for inline error display.
  \item \service{propertyImageUrl} converts the service DTO's image representation
  into a browser-usable URL and supplies the public placeholder when absent.
\end{itemize}

\subsection{UI-to-service flow map}
\begin{tabularx}{\textwidth}{>{\bfseries}p{30mm} p{47mm} X}
\toprule
User action & Angular entry point & Backend operation\\
\midrule
Search homes & Search component / PropertyService & Public Property search via gateway.\\
Register & Register component / AuthService & Identity register; internal User profile create.\\
Open location selector & Registration/seller + LocationService & User states and cities.\\
Create/edit listing & Seller dashboard / PropertyService & Seller-protected property POST/PUT.\\
Upload photos & Seller dashboard / PropertyService & Seller multipart upload; reload seller lists.\\
Verify/deactivate & Admin dashboard / PropertyService & Admin moderation routes.\\
Contact seller & Property details / PropertyService & Buyer contact route; Property asks User.\\
Save/remove & Search/details/cart / PropertyService & Buyer-protected Cart POST/GET/DELETE.\\
\bottomrule
\end{tabularx}

\subsection{Session lifecycle}
After login, \code{AuthService} stores JSON session under \code{ehs_user} and the
JWT under \code{ehs_token}; the signal is updated so the navbar and guards can react.
On application initialization it reads local storage only in a browser platform.
Logout removes both keys and navigates to login. This is a bearer-token SPA pattern;
storage choice is a security trade-off, and token revocation/refresh is not implemented.

\pagebreakdoc
\section{Data Contracts, Validation, and Error Semantics}
\subsection{Property response shape}
Property JSON contains \field{propertyId}, \field{propertyName},
\field{propertyType}, \field{propertyOption}, \field{description}, \field{address},
\field{region}, \field{stateId}, \field{cityId}, \field{priceRange},
\field{initialDeposit}, \field{landmark}, \field{isActive},
\field{deactivationReason}, \field{status}, \field{sellerId}, and \field{images}.
Each image contains \field{imageId}, Base64 \field{image}, and \field{contentType}.
The cart embeds this property DTO in \field{property}; the field is nullable when the
referenced listing is unavailable.

\subsection{Status code guidance}
\begin{longtable}{p{20mm} p{58mm} p{72mm}}
\toprule
Status & Meaning in the current API & Typical source\\
\midrule
\endhead
200 & Query or upload succeeded & Search, detail, location lists, image upload.\\
201 & New resource created & Account registration or property/cart creation (response shape varies).\\
204 & Mutation succeeded without body & Property update, remove cart item, profile compensation.\\
400 & Malformed/invalid business request & Validation, invalid compare pair, image constraints.\\
401 & Authentication or internal-key required/invalid & Invalid login, missing role token, internal route access.\\
403 & Authenticated role lacks permission & ASP.NET Core authorization policy outcome.\\
404 & Requested resource is not available & Property/profile/cart item lookup.\\
409 & Duplicate data conflict & Username/email or duplicate saved property.\\
502 & Downstream service unavailable/failure & Identity profile workflow or Cart property dependency.\\
500 & Unexpected server error & Service exception middleware.\\
\bottomrule
\end{longtable}
Exact serialization can vary by the framework layer; API callers should display safe
messages, not rely on exception text or internal stack details.

\subsection{Validation boundaries}
Angular reactive forms provide immediate field validation, but clients are not
trusted. Backend annotations and explicit controller checks remain authoritative.
Seller property IDs are constrained by owner ID from the token. Cart IDs are
constrained by buyer ID from the token. The frontend's route guard is not a substitute
for either rule. Internal keys are checked only in User Service internal routes.

\section{Security, Logging, and Failure Handling}
\subsection{Security controls present in code}
\begin{itemize}
  \item Password hashes use ASP.NET Core's password hasher.
  \item JWT validation covers issuer, audience, signing key, and lifetime; role-based
  authorization is attached to protected controllers/actions.
  \item Seller ownership and buyer cart ownership are checked server-side.
  \item User profile/contact/internal location routes require a shared internal API key.
  \item CORS origins are configured as \code{http://localhost:4200} for local runs.
  \item Property uploads constrain count, per-file length, and declared MIME type.
\end{itemize}

\subsection{Security gaps / production hardening items}
Development secrets and default admin credentials are present in local configuration
and must not be used in production. Configure secret rotation, TLS, secure cookie
or other token-storage strategy, strict request and upload limits at every hop,
content-signature checks/re-encoding for images, rate limiting, and security headers.
The internal API key is a shared secret and should be rotated and restricted to
service network paths. Review personal contact-data access and audit events for
privacy requirements. No refresh token, token revocation, MFA, or policy-based
fine-grained access model is implemented.

\subsection{Error and dependency behavior}
Identity middleware translates argument errors to 400, selected invalid operations
to 502, and unexpected errors to 500. Cart middleware translates
\code{HttpRequestException} to 502, unauthorized access to 401, and other failures
to 500. User and Property middleware log unexpected exceptions and return a generic
500. Property's User Directory Client maps unresolved city selection to a user-facing
argument error and missing contact to not-found. The gateway streams downstream
status and response content; callers should preserve useful status semantics.

No retry/circuit-breaker policy is configured in the inspected HTTP clients. Blind
retries on non-idempotent operations can create duplicates or ambiguous outcomes;
add retry logic only with idempotency and appropriate timeout policies.

\pagebreakdoc
\section{Startup, Database Initialization, and Operations}
\subsection{Startup behavior}
Each service registers its DbContext, SQL Server connection string, CORS, controllers,
and (where applicable) JWT authentication and authorization. Services expose
development Swagger. On startup, each creates a DI scope and calls
\code{Database.EnsureCreatedAsync()}. User Service seeds its initial location
catalog; Identity Service seeds the admin if no admin account exists. There are no
EF migration bundles or startup schema version checks in the current code.

\subsection{Local startup procedure}
From the solution root, use separate PowerShell terminals:
\begin{verbatim}
dotnet run --project backend\services\UserService\UserService.csproj --launch-profile http
dotnet run --project backend\services\IdentityService\IdentityService.csproj --launch-profile http
dotnet run --project backend\services\PropertyService\PropertyService.csproj --launch-profile http
dotnet run --project backend\services\CartService\CartService.csproj --launch-profile http
dotnet run --project backend\services\ApiGateway\ApiGateway.csproj --launch-profile http
\end{verbatim}
Then start the browser application from \code{frontend} with \code{npm start}. SQL
Server LocalDB and .NET 8 SDK must be installed. The frontend supports an optimized
build via \code{npm run build}; unit tests use \code{npm test -- --watch=false}.

\subsection{Health and diagnostics}
The gateway currently has no dedicated health endpoint. Service liveness can be
checked by contacting each host and its Swagger/API surface in Development; Swagger
availability is not a production health check. Logs use ASP.NET Core logging defaults
with Microsoft framework logs reduced to Warning in the JSON configuration. Add
correlation IDs, structured business-operation logs, redaction, health/readiness
checks, metrics, and central collection before operating multiple replicas.

\section{Test Strategy and Extension Points}
\subsection{Tests in the current repository}
The Angular project contains a small Vitest suite: application rendering,
registration validation, and image URL construction. The backend projects currently
have no checked-in dedicated unit, integration, or API test projects. The backend
directory layout intentionally leaves room for a service-local \code{tests/}
directory as API tests are introduced; no test folder is claimed to exist yet.

\subsection{Recommended test layers}
\begin{enumerate}
  \item \textbf{Unit:} property lifecycle transitions, request validation, claim parsing,
  registration compensation decisions, image URL classification.
  \item \textbf{Service integration:} controller + EF Core against a disposable SQL
  database; verify unique constraints, image rows, state/city constraints, and owner
  isolation.
  \item \textbf{HTTP contract:} gateway route mapping, JWT role behavior, response
  status and JSON casing, multipart boundary and image MIME payload.
  \item \textbf{End-to-end:} register seller, create listing, upload photo, admin
  verify, search as visitor, inspect image, save as buyer, remove saved item.
  \item \textbf{Failure:} downstream timeout/unavailability, oversized request,
  malformed image bytes, invalid token/role, duplicate request and cancellation.
\end{enumerate}

\subsection{API test checkpoints}
\begin{tabularx}{\textwidth}{>{\bfseries}p{34mm} X}
\toprule
Scenario & Measurable expected result\\
\midrule
Image upload & 200, image count persisted, matching content type, returned listing includes Base64 payload.\\
Unauthorized upload & 401/403; no image row is inserted.\\
Cross-seller update/upload & Target property cannot be modified by a different seller.\\
Pending listing search & Pending or inactive property is absent from public search/detail.\\
Verification & Verified and active property appears in public search.\\
Cart duplicate & Second save for same buyer/property returns conflict and creates no second row.\\
Wrong city/state & Property submission is rejected before listing is persisted.\\
Image UI & Browser image source is valid data URL and decoded dimensions are non-zero.\\
\bottomrule
\end{tabularx}

\section{Known Limitations and Future Design Work}
\begin{enumerate}
  \item \textbf{Image transport/storage:} SQL BLOB + Base64 response is easy to reason
  about locally but causes response amplification. Introduce object storage and
  thumbnails once usage warrants it.
  \item \textbf{Search scale:} Add pagination, bounded image metadata, indexes based
  on query plans, and explicit search limits before catalog growth.
  \item \textbf{Cart fan-out:} Replace per-item property fetches with a batch lookup
  or a deliberate cache policy if observed latency is unacceptable.
  \item \textbf{Schema evolution:} Introduce EF Core migrations with backups and
  deployment ordering; do not rely on \code{EnsureCreatedAsync()} for evolving
  long-lived production databases.
  \item \textbf{API resilience:} Add explicit timeouts, resilience policies, request
  size limits, and downstream health/readiness reporting. Preserve idempotency.
  \item \textbf{API contract governance:} Add generated OpenAPI clients or contract
  tests for the manually duplicated property DTO used by Cart Service.
  \item \textbf{Operational security:} Move secrets out of appsettings, rotate
  internal keys, protect the default admin, enforce HTTPS, and define audit/retention
  rules for seller and buyer personal data.
\end{enumerate}

\section{Glossary}
\begin{tabularx}{\textwidth}{>{\bfseries}p{31mm} X}
\toprule
Term & Definition\\
\midrule
Buyer & Consumer role for property discovery, contact access, compare, and saved homes.\\
Seller & Owner role for submitting, editing, and uploading media to own listings.\\
Admin & Moderation role for verifying and deactivating property listings.\\
Profile ID & User Service \field{UserId}, carried as the JWT \field{profileId} claim.\\
Gateway & Browser-facing HTTP path router; forwards requests to the responsible service.\\
Internal API key & Shared service-to-service header used to protect User Service operations.\\
Pending listing & Seller-submitted listing not yet visible to public search.\\
Base64 image DTO & Image bytes represented as Base64 string plus declared MIME type.\\
\bottomrule
\end{tabularx}

\vfill
\begin{center}\small\color{ehgray}End of Low-Level Design\end{center}
\end{document}
